\section{Conclusion}
\label{sec:conclusion}

We introduced YogaAlign, a geometry-grounded extension of Yoga-82, and YogaLLaVA V3, a vision--language model that integrates structured body measurements with continuous angle-aware generation.
YogaAlign combines pose-level domain knowledge with image-specific QA derived from reconstruction-based evidence paths.
YogaLLaVA V3 represents these measurements as metric-aware geometry tokens and predicts continuous values at angle slots.
The experiments demonstrate reliable question-conditioned selection and expression of supplied geometry under both teacher-forced and autoregressive evaluation.
They also define the present scope.
The evaluation covers the geometry-grounded set alone.
YogaLLaVA V3 reasons over monocular reconstruction-derived measurements and is not a substitute for independently validated biomechanical assessment.
Future work should evaluate the domain-knowledge set, modality-matched ablations, human-verified geometric references, pose-held-out generalization, and the consistency of generated corrections with expert instruction.
